% Pursuing Stacks (cotes 134-2 to 134-8): archive cote 134-7, site page 85 (PDF page 86), typescript page 527.
% typescript. Draft compared: ps6.tex:4277-4338.
% AI-transcribed by Claude Opus 5.5 (claude-opus-5-5, adaptive thinking, effort high), Message Batches API, batch ps4 (msgbatch_012Aoc1NuskATPsyBeKSYK7J), 2026-11-12; not yet checked by eye.
\scanpage{134-7/85}
\mspage{527}

ces should be useful in order to give a more or less explicit constructi-%
on of $\tilde{S}$(L.), for given L., in terms of the ``spheres over k'' \xout{\unsure{of}} $S(\underline{\unsure{K}},n)$.

\uline{Remarks}. 1)The superscript 0 for the loop functor $\hand{\Omega}^{0}$ or $\hand{\Omega}^{0}_{ab}$ should
not be confused with the subscript n in the iterated loop-space functor
$\hand{\Omega}_n$ , it is \xout{\unsure{puixis}} added here to suggest that we are taking the \uline{neutral}
\uline{connected component} of the ``true'' full loop-space, \deleted{as} \inserted{this being} imposed by the
restriction of working throughout with 0-connected objects\hand{;} \deleted{--} likewise,
the ``homotopy fiber'' operation in the present context should be viewed
as meaning ``connected component at the marked point of the full homotopy
fiber''. These necessary readjustments of the usual notions is \inserted{being} felt as an
unwelcome feature (of which I have become aware only at this very moment,
I confess, through the writing of the notes). Intuitively, the restric-%
tion to 0-connected \inserted[typed]{(pointed)} homotopy types does'nt seem so imperative, technical
-ly however, when working with \inserted[typed]{semisimplicial} schematic models, we had felt like intro-%
ducing the condition $X_0$=e , which may be viewed as a strong form of a
0-connectedness condition. Truth to tell, it is'nt so clear that this
condition is \xout{of} going to be of great utility - anyhow, it is'nt enough
to ensure what we're really after at present (namely abelian $\hand{\pi}_1$ and tri-%
vial action on the higher $\hand{\pi}_i$'s ), a condition which is anyhow indepen\xout{\illegible{1}}
dant of any 0-connectedness type assumption, and \inserted{is} moreover (it would
seem) stable under the basic fiber and cofiber operations. To sum up,
it may well turn out that we better replace the categories of 0-connected
schematic homotopy types Hotab$_0$(k) and Hot$_0$(k), by slightly larger ones,
so as to get rid of the 0-connectedness restriction. This however is at
present a relatively minor point, and therefore we'll leave the \xout{\unsure{notati-}} \inserted{notations}
ons as they are.

2) Behaviour of L$\widetilde{H}$. with respect to \uline{fibration} sequences \xout{\unsure{is a}}
(instead of cofibration sequences) \xout{are} is \inserted{a} relatively delicate \inserted{matter}, it is
governed by the Leray spectral sequence, whose initial term is the
homology of the base Y with ``coefficients'' in the homology of the \xout{\unsure{(varia-}}
\xout{\unsure{ble)}} fiber. I wonder if there is anything similar for the behaviour
of L$\pi$. with respect to cofibration sequences ?

3) In the display of the main expected properties of the four
basic functors, there is a striking \uline{symmetry}, which we tried to stress
\inserted[typed]{by} \xout{through} the way of presenting the main formulae. One way to express
this symmetry is to say that things look \xout{\illegible{2}} as if there was an auto-%
duality in the pair of categories (H\unsure{o}tab$_0$(k), Hot$_0$(k)), namely a pair
of contravariant involutive functors ($D_{ab}$,D) , each of which interchan\unsure{g}

\margin{85}
